\documentclass[11pt]{article}
\usepackage[margin=0.95in]{geometry}
\usepackage{amsmath,amssymb}
\usepackage{booktabs}
\usepackage{graphicx}
\usepackage[colorlinks=true,linkcolor=blue,urlcolor=blue]{hyperref}
\usepackage{parskip}
\setlength{\parindent}{0pt}

\title{\vspace{-2em}$J/\psi$ radiative corrections:\\
comparison of the two calculations}
\author{V. Kubarovsky}
\date{17 September 2026}

\begin{document}
\maketitle
\vspace{-2.5em}

\noindent
``NLO paper'' is \emph{Fully accounted NLO electromagnetic radiative corrections
to the width of $J/\psi$ decay}; ``the note'' is \emph{Radiative Corrections to
the $J/\psi$ photoproduction reaction}; Ref.~[10] of the NLO paper is Ehlotzky
\& Mitter, Nuovo Cimento \textbf{55A} (1968) 181, abbreviated E\&M.

\section{The closed-form total correction is missing from the NLO paper}

The NLO paper gives $\delta_{self}$~(8), $\delta_{vert}$~(12),
$\delta_{soft}$~(13), $\delta_{hard}$ only as an unevaluated integral~(14), and
the exact Bardin--Shumeiko result~(28). There is no compact closed form for the
total ultrarelativistic correction --- the quantity an experiment applies. It is
E\&M Eq.~(26) with the missing term restored:
\begin{equation*}
\boxed{\;
\delta(m_V,\omega_{cut})=-\frac{4\alpha}{\pi}
\left[
\left(\frac12-\ln\frac{m_V}{m_\ell}\right)\left(\ln r+\frac13\right)
+\frac12\left(L_2(r)-\frac{\pi^{2}}{6}\right)
+\frac{17}{72}+R(r)
+\underbrace{\frac{r^{2}}{16}-\frac{3}{8}r}_{\text{missing from [10]}}
\right]\;}
\end{equation*}
\begin{align*}
r&=\frac{2\omega_{cut}}{m_V}, &
R(r)&=\left(r-\frac{r^{2}}{4}-\frac34\right)\left[\ln\!\left(\frac{m_V}{m_\ell}\sqrt{1-r}\right)-\frac12\right], &
L_2(x)&=-\int_0^x\frac{\ln(1-t)}{t}dt
\end{align*}
For muons add the electron-loop vacuum polarization, $=\delta_{self}$ of Eq.~(8)
with $m_e$. The missing piece is equivalently $+(\alpha/\pi)(3r/2-r^{2}/4)$; it
vanishes as $r\to0$ and reaches $5\alpha/4\pi=0.0029$ at $r=1$.

\textbf{Status of this formula.} It is not exact in the lepton mass, and not the
$m_\ell=0$ limit either: the mass is kept in the logarithms while terms of
relative order $(m_\ell/E)^{2}$, $E=m_V/2$, are dropped. E\&M state this above
their Eq.~(22). Since $(m_e/E)^{2}=1.1\times10^{-7}$ and
$(m_\mu/E)^{2}=4.7\times10^{-3}$, at $r=1$ it is off by $1\times10^{-7}$ for
electrons and $-7\times10^{-5}$ for muons against the exact Eq.~(28).

\textbf{Where it should go: Section~IV}, immediately after Eq.~(14) --- the closed
form \emph{is} the evaluation of that integral combined with Eqs.~(8), (12)
and~(13). Section~IV opens by promising ``compact expressions for all
contributions to decay width''; compact expressions follow for every individual
contribution but not for the sum, which is the only physical combination. Eq.~(16),
$\delta_{FSIC}=3\alpha/4\pi$, then becomes a corollary --- the $r\to1$ limit less
$\delta_{self}$ --- and Figure~4 acquires the function it plots. Since Section~IV
also states that the URA is ``applicable for $\ell=e$'', the electron expression
should be the working formula, with the muon one alongside, its
$(m_\mu/E)^{2}$ error flagged and Eq.~(28) named as the fallback.

\section{The missing term originates in Ref.~[10]}

E\&M Eq.~(23), the differential spectrum, is correct and agrees with Eq.~(14) of
the NLO paper. The term was lost when that spectrum was integrated analytically
in 1968, so Eq.~(24) and Eq.~(26) of [10] are both incomplete. The note
transcribed those equations faithfully and inherited the omission; it has now
been corrected (Section~6).

\textbf{Where to record this:} in Section~IV, immediately after the closed-form
expressions proposed above. Those expressions \emph{are} E\&M Eq.~(26) with the
term restored, so one paragraph placed there attributes the formula, records the
correction, and explains the numerical comparison with [10] all at once.
Suggested wording:

\begin{quote}
The expressions above are Eq.~(26) of Ref.~[10] with the term
$(\alpha/\pi)(3r_{max}/2-r_{max}^{2}/4)$ restored. The differential spectrum,
Eq.~(23) of [10], is correct and agrees with our Eq.~(14); the term was lost in
its analytic integration, so that Eq.~(24), Eq.~(26) and Table~I of [10] are all
low by that amount --- $5\alpha/4\pi=0.0029$ at $r_{max}=1$, and below the
precision quoted in [10] for $r_{max}\lesssim0.1$.
\end{quote}

\section{Table~I of Ref.~[10] was computed with the incomplete formula}

This matters for the validation in Section~IV, because Table~I of [10] is not a
correct benchmark. It was computed from E\&M's own Eq.~(26), and therefore
carries the same omission.

The evidence is not one line but the whole table. Taking the $\delta$ column at
$r_{max}=1.00$, all fourteen mass rows, both channels --- 28 entries:

\begin{center}
\begin{tabular}{lcc}
\toprule
agreement with the $\delta$ printed in Table~I of [10] & electrons & muons\\
\midrule
E\&M Eq.~(26) \emph{as published} (term absent) & 12/14 rows exact & 12/14 rows exact\\
E\&M Eq.~(26) \emph{corrected}                  & 0/14             & 0/14\\
\bottomrule
\end{tabular}
\end{center}

The two misses in each channel are last-digit rounding, not disagreement; the
corrected formula misses every row by $\simeq0.003$. So Table~I of [10] is low by
$(\alpha/\pi)(3r_{max}/2-r_{max}^{2}/4)$ throughout. The agreement reported in
the NLO draft at $\omega/m_e=0.01$ holds only in the small-$r_{cut}$ rows, where
the missing term falls below the three decimals quoted in [10].

Two ways to handle this:

\textbf{(1) State it in the NLO paper.} The paragraph suggested in Section~2
above already covers Table~I of [10]. The sentence claiming the comparison
``succeeds'' should then become a pointer to it, since as written it asserts
agreement with a table that is itself in error.

\textbf{(2) Recompute Table~I at $m_V=M(J/\psi)$.} Since [10] can no longer serve
as the validation target, a table at the physical mass is more useful: it
demonstrates the $\omega$-independence just as well, and the numbers are ones a
$J/\psi$ experiment can use directly. For $m_V=3.0969$~GeV:

\begin{center}
\begin{tabular}{cc rrr c rrr}
\toprule
& & \multicolumn{3}{c}{$e^-e^+$} & & \multicolumn{3}{c}{$\mu^-\mu^+$}\\
\cmidrule(lr){3-5}\cmidrule(lr){7-9}
$r_{cut}$ & $\lg(\omega/m_\ell)$ & $\delta_R$ & $\delta_H$ & $\delta$ & &
$\delta_R$ & $\delta_H$ & $\delta$\\
\midrule
      & $-6$ & $-1.577194$ & $1.481750$ & $-0.095444$ & & $-0.382419$ & $0.376596$ & $-0.005823$\\
0.10  & $-4$ & $-1.225925$ & $1.130481$ & $-0.095444$ & & $-0.259274$ & $0.253452$ & $-0.005823$\\
      & $-2$ & $-0.874655$ & $0.779212$ & $-0.095443$ & & $-0.136130$ & $0.130328$ & $-0.005801$\\
\midrule
      & $-6$ & $-1.577194$ & $1.540558$ & $-0.036636$ & & $-0.382419$ & $0.396794$ & $+0.014375$\\
0.25  & $-4$ & $-1.225925$ & $1.189288$ & $-0.036636$ & & $-0.259274$ & $0.273649$ & $+0.014375$\\
      & $-2$ & $-0.874655$ & $0.838019$ & $-0.036636$ & & $-0.136130$ & $0.150526$ & $+0.014396$\\
\midrule
      & $-6$ & $-1.577194$ & $1.576919$ & $-0.000275$ & & $-0.382419$ & $0.408881$ & $+0.026462$\\
0.50  & $-4$ & $-1.225925$ & $1.225649$ & $-0.000275$ & & $-0.259274$ & $0.285736$ & $+0.026462$\\
      & $-2$ & $-0.874655$ & $0.874380$ & $-0.000275$ & & $-0.136130$ & $0.162613$ & $+0.026483$\\
\midrule
      & $-6$ & $-1.577194$ & $1.603329$ & $+0.026135$ & & $-0.382419$ & $0.416436$ & $+0.034017$\\
1.00  & $-4$ & $-1.225925$ & $1.252060$ & $+0.026135$ & & $-0.259274$ & $0.293291$ & $+0.034017$\\
      & $-2$ & $-0.874655$ & $0.900791$ & $+0.026136$ & & $-0.136130$ & $0.170168$ & $+0.034038$\\
\bottomrule
\end{tabular}
\end{center}
$\delta_R$ for the muon channel includes both self-energy loops. The
$r_{cut}=1.00$ electron entry, $0.026135$, equals
$\delta_{self}+3\alpha/4\pi$ and matches the Bardin--Shumeiko value of Table~III.
The two approaches are not exclusive; (1) is needed in any case.

\section{Figure 4 of the NLO paper and Figure 8 of the note}

These are now the same picture. Figure~8 of the note has been redrawn against
$r_{max}=2k_{max}/M_V$ over the full physical range of each channel, so it
carries the same content as Figure~4, which uses $\omega_{cut}/\omega_{max}$.
The two abscissae are related by $r_{max}=(\omega_{cut}/\omega_{max})\beta^{2}$
with $\beta^{2}=1-4m_\ell^{2}/M_V^{2}$, and the curves agree to machine precision
at every point sampled, in both channels.

\section{Formula comparison: the points that carried a finding}

\begin{center}
\begin{tabular}{llp{8.6cm}}
\toprule
note & NLO & \\
\midrule
(7)        & (6)            & identical (verified symbolically)\\
(11),(12)  & (8)+(12)+(13)  & identical to machine precision\\
(13)       & (14) integrand & identical\\
(14),(15),(19),(20) & ---   & were missing the term; corrected in the note\\
(20)       & ---            & leading logarithm had $m_e$ where $m_\mu$ belongs; corrected\\
(6)        & (1)+(2)        & $4m_\ell^2/m_V$ was not dimensionless; corrected\\
(16)       & (15)           & both take $k_{max}^{\max}=M_V/2$ in the text; the exact endpoint is
                              $M_V/2-2m_\ell^2/M_V$. Corrected in the note; see below\\
---        & (4)            & polarization sum written as for a massless vector\\
---        & (16)           & FSIC $3\alpha/4\pi$; reproduced by the corrected closed form\\
\bottomrule
\end{tabular}
\end{center}

On the last-but-two line: the NLO paper says the cut takes its maximum ``at (15)
or $r_{cut}=r_{max}=1$'', and Table~III is captioned $(r_{max}=1)$, but Eq.~(15)
itself gives $r_{cut}^{\max}=1-4m_\ell^{2}/m_V^{2}=\beta^{2}$, which is $0.99535$
for muons. Taken literally, $\omega_{cut}=m_V/2$ makes the second radicand of
$\rho_1$ in Eq.~(29) negative and $\rho_1$ imaginary, in both channels. The
labels should read $\omega_{cut}=\omega_{max}$, or $r_{cut}=\beta^{2}$.

\section{Status of the note}

Corrected: the missing term in Eqs.~(14), (15), (19), (20); the leading
logarithm of Eq.~(20); the phase-space factor of Eq.~(6); the maximum photon
energy, which now carries the exact endpoint and the bound
$r_{max}\le\beta^{2}$ in Eq.~(16); Fig.~7, recomputed at the $\varepsilon$ stated
in its caption; Fig.~8, redrawn against $r_{max}$ over the full range. A short
paragraph giving the correspondence between the variables of the two papers
($\varepsilon\leftrightarrow\omega$, $k_{max}\leftrightarrow\omega_{cut}$,
$r_{max}\leftrightarrow r_{cut}$, $\ln(M_V/m_\ell)=L/2$) has been added after
Eq.~(17). Equation numbering is unchanged throughout.

Not yet examined: Eqs.~(2)--(5), (17), (18) of the note, and the Bethe--Heitler
formulas behind the $-20\%$ comparison of its Section~VII.

\end{document}
